% !PRINTED 274
\weberpage{274}
%<<CONT
always has two roots. These two roots are distinct from one another
if we exclude substitutions of degree $p$, which indeed cannot occur
in a group whose degree is not divisible by $p$, and which by
\S.\,68, 1.\ are the only ones for which the two roots of (2)
coincide.

These two roots we call the poles of $\Theta$.

We investigate a group $G$ whose degree $n$ is not divisible by $p$,
and which is to be a divisor of $L_p$. If in $G$ the substitutions
\[
\Theta_1,\ \Theta_2,\ \ldots,\ \Theta_{\nu-1},
\]
occur, but no others, which have the same pole $\alpha$, then we
call $\alpha$ \gesperrt{a $\nu$-fold pole} (\S.\,52).

Now if $S$ is any substitution of the group $\mathfrak{E}_{2,p}$, then,
if we transform $L_p$ by $S^{-1}$, we obtain a group $S L_p S^{-1}$
isomorphic with $L_p$, and $G$ passes over, by the same
transformation, into an isomorphic group $S G S^{-1}$.

Although the substitutions of this latter group have no real
coefficients, still each of them has two poles, which one obtains
by applying the substitution $S$ to the poles of $\Theta$. For, if
$\alpha$ is a pole of $\Theta$, we have:
\[
S\,\Theta\,S^{-1}\,S(\alpha) = S\,\Theta(\alpha) = S(\alpha).
\]
We therefore have, as in \S.\,51, 2.:

\gesperrt{1. The group $G$ can be so transformed that an arbitrary
one of its substitutions receives the poles 0 and $\infty$, so that
this substitution becomes multiplicative.}

We now enumerate in order the theorems of \S.\,52, so far as they
come into question here, and shall add an elaboration only there
where the altered hypotheses require it:

\gesperrt{2. If $\alpha$ is a $\nu$-fold pole, the substitutions 1,
$\Theta_1$, $\Theta_2$, . . ., $\Theta_{\nu-1}$ form a cyclic
group. Their substitutions can all (by transformation) be put into
the form:}
\weq{3}{\Theta(\xi) = \gamma^{h\frac{p^2-1}{\nu}}\,\xi, \quad h = 0, 1, \ldots, \nu-1 \quad (\S.\,52, 3.).}
